\documentclass[11pt]{article}

\usepackage[a4paper,margin=1in]{geometry}
\usepackage{amsmath,amssymb,amsthm,mathtools}
\usepackage[T1]{fontenc}
\usepackage{lmodern}
\usepackage{microtype}
\usepackage{hyperref}

\newtheorem{theorem}{Theorem}[section]
\newtheorem{lemma}[theorem]{Lemma}
\newtheorem{proposition}[theorem]{Proposition}
\newtheorem{corollary}[theorem]{Corollary}
\newtheorem{conjecture}[theorem]{Conjecture}
\newtheorem{remark}[theorem]{Remark}
\newtheorem{definition}[theorem]{Definition}

\newcommand{\C}{\mathbb{C}}
\newcommand{\R}{\mathbb{R}}
\newcommand{\eps}{\varepsilon}
\newcommand{\sig}{\sigma}
\newcommand{\conn}{\mathrm{c}}
\newcommand{\tridiag}{\operatorname{tridiag}}
\newcommand{\matA}{\mathbf{A}}
\newcommand{\matD}{\mathbf{D}}
\newcommand{\matI}{\mathbf{I}}

\title{The Connectedness Threshold Problem\\
for the Clean OBC Hatano--Nelson Family}
\author{}
\date{}

\begin{document}
\maketitle

\section{Setup}

We work with the real nonnormal tridiagonal Toeplitz matrices
\[
\matA_n=\tridiag(a,0,b)
=\begin{pmatrix}
0 & b \\
a & 0 & b \\
& \ddots & \ddots & \ddots \\
&& a & 0 & b \\
&&& a & 0
\end{pmatrix}
\in\R^{n\times n},
\quad n>1, \quad 0<a<b.
\]

For a matrix $\matA\in\C^{N\times N}$ and $\eps>0$, its $\eps$-pseudospectrum is
\[
\sig_\eps(\matA)\coloneqq\{z\in\C\colon\|(z\matI-\matA)^{-1}\|>\eps^{-1}\},
\]
where $\|\cdot\|$ is the standard 2-norm.

\begin{lemma}\label{lem:eventual-connectedness}
For every $\eps>0$, there exists an integer $N_\conn$ such that, for all
$n\ge N_\conn$, the pseudospectrum $\sig_\eps(\matA_n)$ is connected.
\end{lemma}

\begin{proof}
Set
\[
q\coloneqq \sqrt{\frac{b}{a}},
\qquad
\matD_n\coloneqq \operatorname{diag}(1,q,\dots,q^{n-1}).
\]
A direct computation gives
\[
\matD_n\matA_n\matD_n^{-1}=\sqrt{ab}\,\tridiag(1,0,1).
\]
Hence $\matA_n$ is similar to the symmetric tridiagonal Toeplitz matrix
$\sqrt{ab}\,\tridiag(1,0,1)$, and therefore its eigenvalues are
\[
\lambda_k^{(n)}=2\sqrt{ab}\cos\Bigl(\frac{k\pi}{n+1}\Bigr),
\qquad k=1,\dots,n.
\]

Now let $\lambda\in\sig(\matA_n)$. If $z\neq \lambda$ and $\mathbf{v}$ is a unit eigenvector
for $\lambda$, then
\[
(z\matI-\matA_n)^{-1}\mathbf{v}=\frac{1}{z-\lambda}\,\mathbf{v},
\]
so
\[
\|(z\matI-\matA_n)^{-1}\|
\ge \frac{1}{|z-\lambda|}.
\]
Thus, with the usual convention that the resolvent norm is $+\infty$ on the
spectrum, we have
\[
B(\lambda,\eps)\subset \sig_\eps(\matA_n),
\qquad
B(\lambda,\eps)\coloneqq\{z\in\C:|z-\lambda|<\eps\}.
\]
In particular,
\[
U_{n,\eps}\coloneqq \bigcup_{k=1}^n B(\lambda_k^{(n)},\eps)
\subset \sig_\eps(\matA_n).
\]

Since the sequence $\lambda_k^{(n)}$ is decreasing in $k$, the mean value theorem
yields
\[
0<\lambda_k^{(n)}-\lambda_{k+1}^{(n)}
=2\sqrt{ab}\left(
\cos\Bigl(\frac{k\pi}{n+1}\Bigr)
-\cos\Bigl(\frac{(k+1)\pi}{n+1}\Bigr)
\right)
\le \frac{2\pi\sqrt{ab}}{n+1}
\]
for $k=1,\dots,n-1$. Choose $N_\conn$ so that
\[
\frac{2\pi\sqrt{ab}}{N_\conn+1}<2\eps;
\]
for instance, one may take
\[
N_\conn=\max\!\left\{2,\Bigl\lfloor \frac{\pi\sqrt{ab}}{\eps}\Bigr\rfloor\right\}.
\]
Then, for every $n\ge N_\conn$,
\[
|\lambda_k^{(n)}-\lambda_{k+1}^{(n)}|<2\eps,
\qquad k=1,\dots,n-1.
\]
Hence consecutive disks $B(\lambda_k^{(n)},\eps)$ and
$B(\lambda_{k+1}^{(n)},\eps)$ overlap, so the union $U_{n,\eps}$ is connected.

It remains to exclude additional connected components of $\sig_\eps(\matA_n)$.
Let $\Omega$ be a connected component of $\sig_\eps(\matA_n)$.
We claim that $\Omega$ contains an eigenvalue of $\matA_n$.
If not, then in fact $\overline{\Omega}\cap\sig(\matA_n)=\varnothing$: indeed, if
$\lambda\in\overline{\Omega}\cap\sig(\matA_n)$, then $\lambda\in\sig_\eps(\matA_n)$
and, since $\sig_\eps(\matA_n)$ is open, a small disk centered at $\lambda$ lies
in $\sig_\eps(\matA_n)$ and meets $\Omega$, contradicting the maximality of the
connected component $\Omega$. Therefore
\[
f(z)\coloneqq \|(z\matI-\matA_n)^{-1}\|
\]
is continuous on $\overline{\Omega}$ and subharmonic on $\Omega$. Also $\Omega$ is
bounded, since $\|(z\matI-\matA_n)^{-1}\|\to 0$ as $|z|\to\infty$.
Moreover, because $\Omega$ is a component of the open set
$\{z:f(z)>\eps^{-1}\}$, we have $f(z)\le \eps^{-1}$ for every
$z\in\partial\Omega$. By the maximum principle for subharmonic functions,
$f(z)\le \eps^{-1}$ throughout $\Omega$, contradicting the definition of
$\Omega\subset\sig_\eps(\matA_n)$. This proves the claim.

Now every connected component of $\sig_\eps(\matA_n)$ contains an eigenvalue, but
all eigenvalues of $\matA_n$ lie in the connected subset
$U_{n,\eps}\subset\sig_\eps(\matA_n)$. Therefore all eigenvalues lie in one and
the same connected component of $\sig_\eps(\matA_n)$, and no other component can
exist. Hence $\sig_\eps(\matA_n)$ is connected for all $n\ge N_\conn$.
\end{proof}

\end{document}